\providecommand{\CoHA}{\mathrm{CoHA}}
\providecommand{\KHA}{\mathrm{KHA}}
\providecommand{\U}{\mathrm{U}}
\providecommand{\Sh}{\mathrm{Sh}}
\providecommand{\C}{\mathbb{C}}
\providecommand{\P}{\mathbb{P}}
\providecommand{\N}{\mathbb{N}}
\providecommand{\bY}{\mathbb{Y}}
\providecommand{\cO}{\mathcal{O}}
\providecommand{\cZ}{\mathcal{Z}}
\providecommand{\Rep}{\mathrm{Rep}}
\providecommand{\Tot}{\mathrm{Tot}}
\providecommand{\tr}{\operatorname{tr}}
\providecommand{\fgl}{\mathfrak{gl}}
\providecommand{\fsl}{\mathfrak{sl}}
\providecommand{\Eone}{E_1}
\providecommand{\Etwo}{E_2}
\providecommand{\kch}{\kappa_{\mathrm{ch}}}
\providecommand{\kcat}{\kappa_{\mathrm{cat}}}
\providecommand{\kBKM}{\kappa_{\mathrm{BKM}}}
\providecommand{\kfib}{\kappa_{\mathrm{fiber}}}

\section*{Conifold Positive Half and the Quantum Toroidal Boundary}
\label{sec:conifold-quantum-toroidal-boundary}

\paragraph{Boundary question.}
Let
\[
  \bY =
  \Tot\bigl(\cO_{\P^1}(-1)\oplus\cO_{\P^1}(-1)\bigr)
\]
be the resolved conifold.  Its non-commutative crepant resolution is the
Klebanov-Witten Jacobi algebra with vertices \(0,1\), arrows
\[
  a_1,a_2:0\to 1,\qquad b_1,b_2:1\to 0,
\]
and potential
\[
  W_{\mathrm{KW}}
  =
  \tr(a_1b_1a_2b_2-a_1b_2a_2b_1).
\]
For a dimension vector \(\gamma=(\gamma_0,\gamma_1)\), set
\[
  \Rep_\gamma(Q_{\mathrm{KW}})
  =
  \bigoplus_{a\in Q_{\mathrm{KW},1}}
  \operatorname{Hom}
  \bigl(\C^{\gamma_{s(a)}},\C^{\gamma_{t(a)}}\bigr),
  \qquad
  G_\gamma=\operatorname{GL}_{\gamma_0}\times
  \operatorname{GL}_{\gamma_1}.
\]
The critical cohomological Hall algebra is
\[
  \CoHA_{\mathrm{con}}
  =
  \bigoplus_{\gamma\in\N^2}
  H^{\bullet}_{c,G_\gamma}
  \bigl(\Rep_\gamma(Q_{\mathrm{KW}}),
  \varphi_{\tr(W_{\mathrm{KW}})_\gamma}\bigr)^\vee ,
\]
with Hall product given by the vanishing-cycle correspondence for short
exact sequences of \(J(Q_{\mathrm{KW}},W_{\mathrm{KW}})\)-modules.
This is the cohomological, rational, positive-half object.  A formula
identifying it directly with
\(\U^+_{q_1,q_2}(\widehat{\widehat{\fgl}}_2)\) conflates this object with
its trigonometric \(K\)-theoretic boundary.

\paragraph{Native operadic level.}
The conifold is a Calabi-Yau threefold, hence the curve output of the
CY-to-chiral construction is native \(\Eone\)-chiral.  The Hall product
above is associative and ordered.  The braided \(\Etwo\) layer belongs to
the Drinfeld centre
\[
  \cZ\bigl(\Rep^{\Eone}(\CoHA_{\mathrm{con}})\bigr)
\]
or to the completed Hall-Drinfeld double after a non-degenerate Hall
pairing has been fixed.  This is the same separation as on an affine
chart:
\[
  \CoHA(\C^3)\simeq Y^+(\widehat{\fgl}_1),
\]
the positive half of the affine Yangian of Schiffmann-Vasserot and
Tsymbaliuk.  The full affine Yangian is the Drinfeld double of
\(Y^+\).  The \(\mathcal W_{1+\infty}\) vacuum module appears through a
completed evaluation action of the full algebra, after doubling and
completion; it is a vertex-algebra representation layer, while the CoHA
itself remains the positive Hall half.

\paragraph{Cohomological shuffle kernel.}
Let \(\epsilon_1,\epsilon_2\) be the two fibre weights of
\(\cO(-1)\oplus\cO(-1)\).  Serre duality supplies the obstruction weight
\(-\epsilon_1-\epsilon_2\).  In additive variables the localized
off-diagonal conifold shuffle factor in the \(0\to 1\) convention is
\[
  \omega_{01}(u)
  =
  \frac{(u+\epsilon_1)(u+\epsilon_2)}
       {u(u+\epsilon_1+\epsilon_2)} .
\]
The numerator records the two arrows \(a_1,a_2:0\to 1\).  The pole at
\(u=0\) is the Hall correspondence pole.  The pole at
\(u=-\epsilon_1-\epsilon_2\) is the Serre-dual obstruction supplied by
the Calabi-Yau threefold potential.  Reversing orientation replaces the
off-diagonal factor by the corresponding \(1\to 0\) kernel.

At the dimension vector \((1,1)\), the potential vanishes on scalar
representations because the four scalar arrow variables commute.  The
Jacobi locus is therefore the whole affine space
\(\C^4_{a_1,a_2,b_1,b_2}\), and the localized Hall product is governed
by the same four arrow weights and the Serre-dual denominator above.
This low-degree computation fixes the kernel normalisation.  The global
presentation theorem requires a separate PBW and wheel-condition proof.

\paragraph{Affine Yangian boundary.}
The cohomological conifold algebra is rational.  Its natural
Yangian-type target is the positive half of the affine super-Yangian
attached to the two-node conifold quiver, often written
\[
  Y^+\bigl(\widehat{\fgl}(1|1)\bigr)_{\mathrm{con}}
\]
or, after choosing the purely even \(A_1^{(1)}\) convention,
\[
  Y^+(\widehat{\fsl}_2).
\]
The choice records the parity convention for the quiver BPS algebra.
The Hall proof obligation is a PBW theorem plus a wheel-condition
identification for the Klebanov-Witten kernel.  Rapcak-Soibelman-Yang-Zhao
(\texttt{arXiv:1810.10402}) prove the double-CoHA action and affine
Yangian identification for the spiked-instanton \(\fgl_1\) setting and
state the broader toric Calabi-Yau expectation.  The conifold case must
therefore cite its own two-node kernel calculation, the Davison-Meinhardt
PBW framework, and the Li-Yamazaki quiver-Yangian dictionary.

\paragraph{Quantum toroidal boundary.}
Quantum toroidal algebras are trigonometric double-loop objects.  For
parameters \(q_1q_2q_3=1\), their current relations use multiplicative
kernel factors
\[
  g(z,w)=(z-q_1w)(z-q_2w)(z-q_3w)
\]
and the toroidal Serre relations.  Feigin-Jimbo-Miwa-Mukhin
(\texttt{arXiv:1309.2147}, Adv. Math. 2016) study representation
branching for quantum toroidal \(\fgl_n\).  Negu\c{t}
(\texttt{arXiv:1302.6202}) proves that the quantum toroidal algebra of
\(\fgl_n\) is isomorphic to the double shuffle algebra of Feigin-Odesskii
for the cyclic quiver; his thesis version
(\texttt{arXiv:1504.06525}) gives the K-theoretic cyclic-quiver
variety construction.  Varagnolo-Vasserot
(\texttt{arXiv:2011.01203}) prove K-theoretic Hall algebra
identifications with positive halves of quantum toroidal quantum groups
for affine preprojective type and compare type \(A\) super-toroidal
quantum groups with \(K\)-theoretic Hall algebras of quivers with
potential.

These results justify the following boundary statement.  After replacing
the critical cohomological Hall algebra by the \(K\)-theoretic Hall
algebra and after matching the Klebanov-Witten wheel conditions with
the cyclic \(A_1^{(1)}\) current relations, the trigonometric boundary
is expected to be
\[
  \KHA^+(\bY)
  \simeq
  \U^+_{q_1,q_2}
  \bigl(\widehat{\widehat{\fgl}}_2\bigr)
\]
up to the central Heisenberg extension determined by the \(\fgl_2\)
convention; the \(\fsl_2\) convention is obtained by removing this
central sector.  The displayed formula is a conditional
\(K\)-theoretic shuffle comparison.  Its missing local ingredient is
the explicit wheel-to-current calculation for the Klebanov-Witten
two-node kernel.

\paragraph{Rational limit and Miki scope.}
The rational degeneration from quantum toroidal to affine Yangian is
visible on the structure function.  Put \(q_a=\exp(\hbar h_a)\) with
\(h_1+h_2+h_3=0\), \(x=\exp(\hbar u)\), and
\[
  \varphi_{\mathrm{tor}}(x)
  =
  \prod_{a=1}^{3}
  \frac{1-q_a^{-1}x}{1-q_a x}.
\]
Then
\[
  \lim_{\hbar\to 0}\varphi_{\mathrm{tor}}(\exp(\hbar u))
  =
  \prod_{a=1}^{3}\frac{u-h_a}{u+h_a}
  =
  g_{\mathrm{aff}}(u),
\]
the affine-Yangian structure function used in the \(\C^3\) chart
calculation.  Miki's automorphism is an automorphism of the toroidal
double-loop algebra, exchanging the two loop directions and generating
the \(SL_2(\mathbb Z)\) symmetry of the torus.  In the rational
Yangian limit the second loop has degenerated; only the permutation
symmetry of the parameters \(h_1,h_2,h_3\) remains on the shuffle
kernel.  Thus Miki symmetry is a quantum-toroidal datum, while the
affine Yangian keeps its rational \(S_3\)-symmetric structure function.

\paragraph{Citation separation.}
The citation string
``Feigin-Jimbo-Mukhin-Vilkoviskiy 2021'' is unrelated to the conifold
CoHA comparison; the matching arXiv identifier belongs to Zenkevich's
note on the pentagon identity in the Ding-Iohara-Miki algebra.  The
quantum-toroidal branching reference used here is
Feigin-Jimbo-Miwa-Mukhin \(\texttt{arXiv:1309.2147}\).  The conifold
topological-string reference Awata-Feigin-Shiraishi
\(\texttt{arXiv:1112.6074}\) constructs refined-topological-vertex
matrix elements from Miki's \(\mathcal W_{1+\infty}\) or DIM
intertwiners.  It supplies a refined-vertex amplitude comparison.  It
leaves the critical CoHA construction for \(\bY\) and the
Klebanov-Witten wheel-to-current theorem as external obligations.

\paragraph{Branch separation.}
The conifold toric branch, the \(\C^3\) affine-Yangian chart branch,
and the \(K3\times E\) Hall-Borcherds branch are separate constructions.
For the conifold, the invariant used in this note is the direct McKay
value
\[
  \kch(A_{\bY})=1 .
\]
A Borcherds product is absent from the non-compact conifold argument,
so no \(\kBKM\) value is assigned.  The fibre invariant \(\kfib\) is
likewise absent.  The \(K3\times E\) constants
\[
  \{\kcat,\kch^{\mathrm{Heis}},\kBKM,\kfib\}
  =
  \{0,3,5,24\}
\]
belong to the Hall-Drinfeld double and Borcherds-product branch; they
are outside the conifold quantum-toroidal comparison.  Root-of-unity
truncation for \(U_{q,t}(\widehat{\widehat{\fgl}}_1)\) or
\(U_{q,t}(\widehat{\widehat{\fgl}}_2)\) remains an open quotient problem,
so finite module counts and Verlinde-type specialisation theorems remain
outside the scope of this note.

\paragraph{Proof boundary.}
The proved content retained here is the construction of the
Klebanov-Witten critical CoHA, the additive conifold shuffle kernel, the
\(\Eone\) native level for a CY\(_3\) output, the positive-half
identification \(\CoHA(\C^3)=Y^+(\widehat{\fgl}_1)\) on affine charts,
and the rational limit from toroidal structure functions to affine
Yangian structure functions.  The \(K\)-theoretic equality with
\(\U^+_{q_1,q_2}(\widehat{\widehat{\fgl}}_2)\) is recorded as a
construction-level target: prove the conifold wheel conditions, match
the central Heisenberg extension, and check the PBW/no-extra-relations
ideal after completion.
